% ======================================================================
\chapter{Designs}

\section{System Architecture Overview}
ScamGuard is a three-tier web application. The \textbf{presentation tier} is a Next.js front end hosted on Vercel; it proxies all API calls through its own origin. The \textbf{application tier} is a FastAPI back end hosted on Render, which runs the detector, the intent rules, the cross-domain models, authentication and the safety tools in one process. The \textbf{data tier} is a PostgreSQL database on Neon that stores users, scans, cases and reports. Models are trained offline and stored as versioned files that the back end loads once at start-up. Figure~\ref{fig:arch} shows the components and their connections.

\begin{figure}[H]
\centering
\resizebox{\linewidth}{!}{%
\begin{tikzpicture}[
  box/.style={draw, rounded corners=2pt, align=center, font=\small, minimum width=3.8cm, minimum height=12mm, inner sep=4pt, fill=white},
  off/.style={box, fill=gray!15},
  arr/.style={-{Stealth[length=2mm]}, thick}]
\def\dx{4.6}
% offline row
\node[off] (xnb)   at (0,0)        {Cross-domain notebook\\(UCI + India)};
\node[box] (xdm)   at (\dx,0)      {Cross-domain models\\(baseline, adapted)};
\node[off] (train) at (2*\dx,0)    {Offline training on\\ScamGuard-IN-2026};
\node[off] (reg)   at (3*\dx,0)    {Model registry\\(versioned files)};
% request row
\node[box] (ui)    at (0,-2.2)     {Web app\\(Next.js, Vercel)};
\node[box] (api)   at (\dx,-2.2)   {FastAPI back end\\(Render)};
\node[box] (pre)   at (2*\dx,-2.2) {Cleaner +\\tokenizer};
\node[box] (clf)   at (3*\dx,-2.2) {TF-IDF +\\logistic regression};
% response row
\node[box] (db)    at (0,-4.4)     {PostgreSQL\\(Neon)};
\node[box] (expl)  at (\dx,-4.4)   {Scam X-Ray,\\TrustLens};
\node[box] (fuse)  at (2*\dx,-4.4) {Fusion, guards,\\risk score};
\node[box] (rules) at (3*\dx,-4.4) {Intent and\\protective rules};
% helpers
\node[box] (cop)   at (\dx,-6.6)   {AI copilot\\(Gemini, capped)};
\node[box] (rag)   at (2*\dx,-6.6) {Known-pattern\\retrieval (52 entries)};
\draw[arr] (xnb) -- (xdm); \draw[arr] (train) -- (reg); \draw[arr] (reg) -- (clf);
\draw[arr, <->] (api) -- (xdm);
\draw[arr] (ui) -- (api); \draw[arr] (api) -- (pre); \draw[arr] (pre) -- (clf);
\draw[arr] (clf) -- (rules); \draw[arr] (rules) -- (fuse); \draw[arr] (fuse) -- (expl);
\draw[arr] (expl) -- (db); \draw[arr] (db) -- (ui);
\draw[arr] (rag) -- (expl); \draw[arr] (expl) -- (cop);
\end{tikzpicture}}
\caption{Architecture of the deployed system. Shaded parts run offline.}
\label{fig:arch}
\end{figure}


\section{Module/Component Design}
Table~\ref{tab:subsys} lists the modules of ScamGuard, their responsibility and their interfaces.

\begin{table}[H]
\centering
\caption{Modules of ScamGuard and their interfaces}
\label{tab:subsys}
\begin{tabular}{|L{3.0cm}|L{5.3cm}|L{4.4cm}|}
\hline
\textbf{Module} & \textbf{Responsibility} & \textbf{Interface} \\
\hline
Dataset builder & Collect sources, audit labels, group campaigns, split & Python script; writes one CSV file \\
\hline
Training pipeline & Train and compare models; register the selected one & Model registry (vectorizer and model files) \\
\hline
Detector & Clean, classify, apply rules, fuse, explain & Python function \texttt{predict(text)} \\
\hline
Cross-domain study & Audit and clean UCI and Indian data, run the 7 experiments over 3 seeds, save results and models & Jupyter notebook; CSV tables, figures, \texttt{results.json}, joblib models \\
\hline
Cross-domain service & Serve the baseline and adapted models with explanations & \texttt{POST /public/cross-domain/classify}, \texttt{GET /public/cross-domain/results} \\
\hline
Back end (API) & Authentication, rate limits, storage, reports, safety tools & REST over HTTPS, JSON \\
\hline
Front end & Quick check, analysis, dashboard, tools, cross-domain page, Hindi & Browser; calls the REST API \\
\hline
Streamlit app & Stand-alone demonstration of the study & Loads the saved models; never retrains \\
\hline
Copilot & Answers follow-up questions & REST endpoint; external LLM API with limits \\
\hline
\end{tabular}
\end{table}

\section{Detailed Design}

\subsection{Client--server communication}
The web client and the back end communicate over HTTPS with JSON bodies. Figure~\ref{fig:seq} shows the sequence for one analysis.

\begin{figure}[H]
\centering
\resizebox{0.95\linewidth}{!}{%
\begin{tikzpicture}[font=\small, >={Stealth[length=2.2mm]}, msg/.style={->, thick}, ret/.style={->, thick, dashed}, every node/.append style={inner sep=1.5pt}, lbl/.style={above, fill=white, font=\small}]
\foreach \x/\n/\t in {0/u/User, 3.6/c/Web client\\(Next.js), 7.6/a/FastAPI\\back end, 11.6/m/Detector, 15.2/d/Database} {
  \node[draw, rounded corners=2pt, fill=gray!12, minimum width=2.7cm, minimum height=10mm, align=center] (\n) at (\x,0) {\t};
  \draw[gray] (\x,-0.55) -- (\x,-10.2);
}
\draw[msg] (0,-1.2) -- node[lbl]{1. paste message} (3.6,-1.2);
\draw[msg] (3.6,-2.1) -- node[lbl]{2. POST /analyze} (7.6,-2.1);
\node[draw, fill=white, align=center, font=\footnotesize, minimum width=2.9cm] at (7.6,-3.2) {3. check token, rate\\and size limits; clean};
\draw[msg] (7.6,-4.3) -- node[lbl]{4. predict(text)} (11.6,-4.3);
\node[draw, fill=white, align=center, font=\footnotesize, minimum width=2.5cm] at (11.6,-5.4) {5. classifier, rules,\\fusion, explanation};
\draw[ret] (11.6,-6.5) -- node[lbl]{6. verdict, score, phrases} (7.6,-6.5);
\draw[msg] (7.6,-7.3) -- node[lbl]{7. save scan} (15.2,-7.3);
\draw[ret] (15.2,-8.0) -- node[lbl]{8. saved} (7.6,-8.0);
\draw[ret] (7.6,-8.8) -- node[lbl]{9. 200 OK (JSON)} (3.6,-8.8);
\draw[ret] (3.6,-9.6) -- node[lbl]{10. show verdict} (0,-9.6);
\end{tikzpicture}}
\caption{Sequence diagram of one message analysis}
\label{fig:seq}
\end{figure}

\begin{itemize}
\item \textbf{Endpoints.} \texttt{POST /api/v1/public/quick-check} serves visitors without an account (10 requests per minute); \texttt{POST /api/v1/messages/analyze} serves signed-in users (30 per minute); \texttt{POST /api/v1/public/cross-domain/classify} serves the cross-domain page (30 per minute). Other routes handle authentication, history, cases, analytics, the copilot and the safety tools.
\item \textbf{Authentication.} Signed access tokens with refresh tokens; passwords stored as bcrypt hashes.
\item \textbf{Robustness.} Message text is limited to 4,000 characters (2,000 for the public endpoints).
\end{itemize}

\subsection{Dataset construction}
Algorithm~\ref{alg:build} summarises how ScamGuard-IN-2026 is built and split.

\begin{algorithm}[H]
\caption{Building and splitting ScamGuard-IN-2026}
\label{alg:build}
\begin{algorithmic}[1]
\Require sources R1--R3, synthetic templates S, audit files
\State normalise anonymised entities in R1; mask personal data
\State group real messages into campaigns by MinHash ($J\ge0.6$)
\For{each reported campaign}
  \If{strong fraud cue and no genuine cue} \State label scam \Else \State apply the manual decision (scam, genuine or drop) \EndIf
\EndFor
\State re-read rows that a cross-validated model strongly contradicts
\State assign each campaign to train, development or test by a salted hash
\State add synthetic messages to training only; remove exact test duplicates from training
\State \Return labelled dataset with split column
\end{algorithmic}
\end{algorithm}

\subsection{Cross-domain experiment design}
Algorithm~\ref{alg:xdomain} gives the protocol of the cross-domain study. All choices are made on source data; the final target test set is used only for scoring.

\begin{algorithm}[H]
\caption{Cross-domain evaluation and lightweight domain adaptation}
\label{alg:xdomain}
\begin{algorithmic}[1]
\Require cleaned source set $A$ (UCI), cleaned target set $B$ (Indian), seeds $\{42, 123, 2026\}$, fractions $\{0, 0.1, 0.2, 0.4\}$
\State choose $\alpha$ (Naive Bayes) and $C$ (Logistic Regression) by 5-fold cross-validation on $A_{train}$
\For{each seed}
  \State split $A$ 80/20 into $A_{train}, A_{test}$ and $B$ 70/30 into pool $P$ and test $B_{test}$ (stratified)
  \For{each fraction $f$}
    \State $S_f \gets$ stratified sample of $f\cdot|P|$ messages from $P$
    \State fit TF-IDF and the classifier on $A_{train}\cup S_f$
    \State score $A_{test}$ and $B_{test}$; store predictions
  \EndFor
\EndFor
\State report mean $\pm$ SD over seeds; McNemar test of each adapted model against $f=0$
\end{algorithmic}
\end{algorithm}

\subsection{Inference procedure}
\begin{algorithm}[H]
\caption{Scoring one message in the live service}
\begin{algorithmic}[1]
\Require message $x$; frozen vectorizer and model; intent rules
\State $t\gets$ \textsc{Tokenize}(\textsc{Clean}($x$))
\State $p\gets\sigma(\mathbf{w}^\top\phi(t)+b)$
\State find matched intents $I$ and protective signals $P$
\If{a sentence asks for an OTP after a ``never share'' warning} \State add decoy intent \EndIf
\State $s\gets\sigma(\operatorname{clip}(\operatorname{logit}p,\pm4)+5(E^+-E^-))$
\If{genuine alert or notice with no blocking cue} \State $s\gets\min(s,0.15)$ \EndIf
\State \Return verdict $\mathbf{1}[s\ge0.5]$, $p$, $s$, highlighted phrases, explanation
\end{algorithmic}
\end{algorithm}

\section{Data Design}

\subsection{Database}
Table~\ref{tab:db} lists the main tables of the PostgreSQL database.

\begin{table}[H]
\centering
\caption{Main database tables}
\label{tab:db}
\begin{tabular}{|L{3.6cm}|L{9.9cm}|}
\hline
\textbf{Table} & \textbf{Contents} \\
\hline
\texttt{users} & Account, hashed password, role, preferences \\
\hline
\texttt{refresh\_tokens} & Issued refresh tokens for rotation and revocation \\
\hline
\texttt{messages} & Each analysed input with its type (text, URL, e-mail, image, PDF, QR) \\
\hline
\texttt{predictions} & Verdict, classifier probability, risk score, category, explanation, model version \\
\hline
\texttt{cases}, \texttt{case\_scan\_links}, \texttt{case\_timeline\_entries} & Investigation cases, the scans linked to them and their notes \\
\hline
\texttt{community\_reports} & One report per user for a number, link domain, UPI ID or e-mail \\
\hline
\texttt{alerts}, \texttt{audit\_log}, \texttt{model\_registry\_meta} & Alerts, security audit trail and registered model versions \\
\hline
\end{tabular}
\end{table}

\subsection{Datasets}
\textbf{ScamGuard-IN-2026} has seven attributes per row: \texttt{text}, \texttt{label} (scam or legitimate), \texttt{source}, \texttt{subtype}, \texttt{is\_synthetic}, \texttt{campaign} (near-duplicate group) and \texttt{split} (train, development, test or 2026 test).

The \textbf{cross-domain datasets} are stored after cleaning with the attributes \texttt{message}, \texttt{label} (spam = 1, ham = 0) and \texttt{domain} (UCI or India). The study also writes every result table as CSV, the error-annotation sheet with a \texttt{manual\_category} column, a \texttt{results.json} summary and four joblib models: Naive Bayes and Logistic Regression, each in a baseline and an adapted version.
